Un circuito serie idealizado permanece cerrado y contiene una resistencia \(R=1\,\Omega\) y una inductancia \(L=1\) H. Sus caídas de tensión son \(Ri\) y \(Li'\), y su suma iguala la tensión de la fuente. En cada ciclo de tres segundos, la fuente proporciona \(6\) V durante el primer segundo y \(0\) V durante los dos siguientes; la operación se repite indefinidamente. Los valores son sintéticos.

Formule la entrada periódica y la ecuación para la corriente. Elija y justifique una herramienta de resolución adecuada para obtener la operación periódica. Determine qué corriente debe existir al comienzo del primer ciclo para que la respuesta sea periódica desde ese instante, y dé la corriente durante un ciclo completo. Distinga esa inicialización de lo que ocurre al empezar con otra corriente. Obtenga además la corriente media de un ciclo integrando el balance de tensión, sin calcular por separado la integral de cada tramo de la solución.
